\chapter{Introduction}
\label{chap:1}
Rigid-body dynamics is an old subject that has been rejuvenated and transformed by the computer. Today, we can find dynamics calculations in computer games, in animation and virtual-reality software, in simulators, in motion control systems, and in a variety of engineering design and analysis tools. In every case, a computer is calculating the forces, accelerations, and so on, associated with the motion of a rigid-body approximation of a physical system. The main purpose of this book is to present a collection of efficient algorithms for performing various dynamics calculations on a computer. Each algorithm is described in detail, so that the reader can understand how it works and why it is efficient; and basic concepts are explained, such as recursive formulation, branch-induced sparsity and articulated-body inertia. Rigid-body dynamics is usually expressed using 3D vectors. However, the subject of this book is dynamics algorithms, and this subject is better expressed using 6D vectors. We therefore adopt a 6D notation based on spatial vectors, which is explained in Chapter 2. Compared with 3D vectors, spatial notation greatly reduces the volume of algebra, simplifies the tasks of describing and explaining a dynamics algorithm, and simplifies the process of implementing an algorithm on a computer. This chapter provides a short introduction to the subject matter of this book. It says a little about dynamics algorithms, a little about spatial vectors, and it explains how the book is organized. 1.1 Dynamics Algorithms The dynamics of a rigid-body system is described its equation of motion, which specifies the relationship between the forces acting on the system and the accelerations they produce. A dynamics algorithm is a procedure for calculating the numeric values of quantities that are relevant to the dynamics. We will be concerned mainly with algorithms for two particular calculations:

forward dynamics: the calculation of the acceleration response of a given rigid-body system to a given applied force; and inverse dynamics: the calculation of the force that must be applied to a given rigid-body system in order to produce a given acceleration response. Forward dynamics is used mainly in simulation. Inverse dynamics has a variety of uses, such as: motion control systems, trajectory planning, mechanical design, and as a component in a forward-dynamics calculation. In Chapter 9 we will consider a third kind of calculation, called hybrid dynamics, in which some of the acceleration and force variables are given and the task is to calculate the rest. The equation of motion for a rigid-body system can be written in the following canonical form: τ = H(q) ¨q + C(q, ˙q) . (1.1) In this equation, q, ˙q and ¨q are vectors of position, velocity and acceleration variables, respectively, and τ is a vector of applied forces. H is a matrix of inertia terms, and is written H(q) to show that it is a function of q. C is a vector of force terms that account for the Coriolis and centrifugal forces, gravity, and any other forces acting on the system other than those in τ. It is written C(q, ˙q) to show that it depends on both q and ˙q. Together, H and C are the coefficients of the equation of motion, and τ and ¨q are the variables. Typically, one of the variables is given and the other is unknown. Although it is customary to write H(q) and C(q, ˙q), it would be more accurate to write H(model, q) and C(model, q, ˙q), to indicate that both H and C depend on the particular rigid-body system to which they apply. The symbol model refers to a collection of data that describes a particular rigid-body system in terms of its component parts: the number of bodies and joints, the manner in which they are connected together, and the values of every parameter associated with each component (inertia parameters, geometric parameters, and so on). We call this description a system model in order to distinguish it from a mathematical model. The difference is this: a system model describes the system itself, whereas a mathematical model describes some aspect of its behaviour. Equation 1.1 is a mathematical model. On a computer, model would be a variable of type ‘system model’, and its value would be a data structure containing the system model of a specific rigid-body system. Let us encapsulate the forward and inverse dynamics calculations into a pair of functions, FD and ID. These functions satisfy ¨q = FD(model, q, ˙q, τ) (1.2) and τ = ID(model, q, ˙q, ¨q) . (1.3) On comparing these equations with Eq. 1.1, it is obvious that FD and ID must evaluate to H−1(τ −C) and H¨q +C, respectively. However, the point of these

equations is that they show clearly the inputs and outputs of each calculation. In particular, they show that the system model is an input in both cases. Thus, there is an expectation that the algorithms implementing FD and ID will work for a class of rigid-body systems, and will use the data in the system model to work out the dynamics of the particular rigid-body system described by that model. We shall use the name model-based algorithm to refer to algorithms that work like this. The big advantage of this approach is that a single piece of computer code can be written, tested, documented, and so on, to calculate the dynamics of any rigid-body system in a broad class. The two main classes of interest are called kinematic trees and closed-loop systems. Roughly speaking, a kinematic tree is any rigid-body system that does not contain kinematic loops, and a closedloop system is any rigid-body system that is not a kinematic tree.1 Calculating the dynamics of a kinematic tree is significantly easier than for a closed-loop system. Model-based algorithms can be classified according to their two main attributes: what calculation they perform, and what class of system they apply to. The main body of algorithms in this book consists of forward and inverse dynamics algorithms for kinematic trees and closed-loop systems. 1.2 Spatial Vectors A rigid body in 3D space has six degrees of motion freedom; yet we usually express its dynamics using 3D vectors. Thus, to state the equation of motion for a rigid body, we must actually state two vector equations: f = m aC and nC = I ˙ω + ω × Iω . (1.4) The first expresses the relationship between the force applied to the body and the linear acceleration of its centre of mass. The second expresses the relationship between the moment applied to the body, referred to its centre of mass, and its angular acceleration. In spatial vector notation, we use 6D vectors that combine the linear and angular aspects of rigid-body motion. Thus, linear and angular acceleration are combined to form a spatial acceleration vector, force and moment are combined to form a spatial force vector, and so on. Using spatial notation, the equation of motion for a rigid body can be written f = Ia + v ×∗Iv , (1.5) where f is the spatial force acting on the body, v and a are the body’s spatial velocity and acceleration, and I is the body’s spatial inertia tensor. The symbol ×∗denotes a spatial vector cross product. One obvious feature of this equation 1A precise definition will be given in Chapter 4, including a definition of ‘kinematic loop.’

is that there are some name clashes with Eq. 1.4. We can solve this problem by placing hats above the spatial symbols (e.g. ˆf). Spatial notation offers many more simplifications than just the equation of motion. For example, if two rigid bodies are joined together to form a single rigid body, then the spatial inertia of the new body is given by the simple formula Inew = I1 + I2 , where I1 and I2 are the inertias of the two original bodies. This equation replaces three in the 3D vector approach: one to compute the new mass, one to compute the new centre of mass, and one to compute the new rotational inertia about the new centre of mass. The overall effect of all these simplifications is that spatial notation typically cuts the volume of algebra by at least a factor of 4 compared with standard 3D vector notation. With the barrier of algebra out of the way, the analyst is free to state a problem more succinctly, to solve it in fewer steps, and to arrive at a more compact solution. Another benefit of spatial vectors is that they simplify the process of writing computer code. The resulting code is shorter, and is easier to read, write and debug, but is no less efficient than software using 3D vectors. For example, some programming languages will allow you to implement Eq. 1.5 with a statement like f = I*a + v.cross(I*v); in which the type of each variable has been declared to the compiler. Thus, the compiler knows that I and v denote a rigid-body inertia and a spatial motion vector, respectively, so it can compile the expression I*v into code that calls a routine in the spatial arithmetic library which is designed to perform this particular multiplication efficiently. 1.3 Units and Notation Angles are assumed to be measured in radians. Apart from that, the equations in this book are independent of the choice of units, and require only that the chosen system of units be consistent. In the few places where units are mentioned explicitly, SI units are used. The mathematical notation generally follows ISO guidelines. Thus, variables are set in italic; constants and functions are set in roman; and vectors and matrices are set in bold italic. Vectors are denoted by lower-case letters,2 and matrices by upper-case. A short list of symbols can be found on page 265, and some symbols have entries in the index. A few details are worth mentioning here. The symbols 0 and 1 denote the zero and identity matrices, respectively, and 0 also denotes the zero vector. The superscript −T denotes the transpose of the inverse, so A−T means (A−1)T. An 2The vector C in Eq. 1.1 is the only exception to this rule.

expression like a× denotes an operator that maps b to a × b. If a quantity is described as an array, then it is a numbered list of things. If λ is an array then element i of λ is written λ(i). If a symbol has a leading superscript (e.g. Av) then that superscript identifies a coordinate system. Coordinate systems can also appear in subscripts. Spatial vectors are sometimes marked with a hat, and coordinate vectors with an underline, as explained in Chapter 2. 1.4 Readers’ Guide The contents of this book can be divided into three parts: preparation, main algorithms and additional topics. The preparatory chapters are 2, 3 and 4; and the main algorithms are are described in Chapters 5 to 8. Readers with enough background may find they can skip straight to Chapter 5. Chapters 9 onwards assume a basic understanding of earlier material, but are otherwise self-contained. Chapter 2 describes spatial vector algebra from first principles; and Chapter 3 explains how to formulate and analyse the equations of motion for a rigidbody system. These are the two most mathematical chapters, and they both cover their subject matter in greater depth than the minimum required for the algorithms that follow. Chapter 4 describes the various components of a system model. Many of the quantities used in later chapters are defined here. Chapters 5, 6 and 7 describe the three best known algorithms for kinematic trees: the recursive Newton-Euler algorithm for inverse dynamics, and the composite-rigid-body and articulated-body algorithms for forward dynamics, in that order. Chapter 5 also explains the concept of recursive formulation, which is the reason for the efficiency of these algorithms. Chapter 8 then presents several techniques for calculating the dynamics of closed-loop systems. These four chapters are arranged approximately in order of increasing complexity (i.e., the algorithms get progressively more complicated). Chapters 9, 10 and 11 then tackle a variety of extra topics. Chapter 9 considers several more algorithms, including hybrids of forward and inverse dynamics and algorithms for floating-base systems. Chapter 10 considers the issues of numerical accuracy and computational efficiency; and Chapter 11 examines the dynamics of rigid-body systems that are subject to contacts and impacts. Readers will find examples dotted unevenly throughout the text. In some cases, they illustrate a concept that is covered in the main text. In other cases, they describe ideas that are more easily explained in the form of an example. Readers will also find that every major algorithm is presented both as a set of equations and as a pseudocode program. In many cases, the two are side-byside. The pseudocode is intended to be self-explanatory and easy to translate into real computer code. Many people, including the author, find it helps to have some working source code to hand when studying a new algorithm. No source code is distributed with this book, on the grounds that such software becomes obsolete too quickly.

Readers should instead look for this code on the Web. At the time of writing, the author’s web site3 contained source code for most of the algorithms described here, but readers may also be able to find useful software elsewhere. 1.5 Further Reading A variety of books on rigid-body dynamics have been published in recent years, although many are now out of print. The following are suitable introductory texts: Amirouche (2006), Moon (1998), Shabana (2001) and Huston (1990). The first three are formal teaching texts with copious examples and homework exercises, while Huston’s book is more tutorial in nature. All are designed to be accessible to undergraduates. More advanced texts include Stejskal and Val´aˇsek (1996), Roberson and Schwertassek (1988) and Wittenburg (1977). For a different perspective, try Coutinho (2001), which is aimed more at computer games and virtual reality. Much of the material in this book has its origins in the robotics community. Other books on robot dynamics include Balafoutis and Patel (1991), Lilly (1993), Yamane (2004) and the author’s earlier work, Featherstone (1987), which has been superseded by the present volume. Sections and chapters on dynamics can be found in Angeles (2003), Khalil and Dombre (2002), Siciliano and Khatib (2008) and various other robotics books. Readers with an interest in 6D vectors of various kinds may find the following books and articles of interest: Ball (1900), Brand (1953), von Mises (1924a,b), Murray et al. (1994) and Selig (1996). Also, the treatment of spatial vectors in Featherstone (1987) is a little different from that presented here. 3The URL is (or was) http://users.rsise.anu.edu.au/\textasciitilde{}roy/spatial/. If that doesn’t work, then try searching for ‘Roy Featherstone’.
